\section{The machine}
\label{sec:machine}

\begin{question}
  \textbf{What is different about this machine, compared with its description?} Which nodes can be used, what limits them, and what
  does that do to every number in this report?
\end{question}

Of the 52 nodes, 43 could be used and 40 were clean until November~6, the morning the administrators cleared a 200\,W power limit from six nodes; now 48 can be used and 45 are clean. A capped GPU is about
\TriCapFactor{} times slower, and one node is still capped. One cleared node now fails an NCCL test, and two others are out. CPU clocks fall under high GPU power on all but one node, so \textbf{every Qwen step
time in this report is throttled, except under pipeline layouts}. Three nodes spend about twice as long at low CPU clocks and slow any job that holds them. A kernel change during the
campaign had no measurable effect, so results before and after it are compared directly. \Cref{fig:machine_status} shows the nodes.

\begin{figure}[h]
  \centering\fig{machine_status}
  \caption{Top: CPU clock under full GPU power in the health check, per node, with load on 72 of 96 cores. Bottom: share of CPU samples below 500\,MHz per
    node and run, in 13-node Qwen FSDP training at 12k tokens per GPU, from the node telemetry. Red: GPUs capped at 200\,W for part of the campaign; gray: node out.}
  \label{fig:machine_status}
\end{figure}

\paragraph{Nodes.} \texttt{trig0146} is out: it shows 7 of its 8 GPUs. \texttt{trig0122} is in maintenance and excluded until the administrators confirm a firmware update
(its \texttt{mlx5\_1} runs 28.35.2000, other nodes 28.47.1026). That leaves 50. \texttt{trig0145}, still capped and reserved by the administrators, and \texttt{trig0140}, which fails an NCCL test, leave 48, and the 3 slow ones leave 45 clean nodes. While the caps lasted those numbers were 43 and 40, so the runs so far stop at 39--43 nodes.
\texttt{trig0122} also read 200\,W when last read, before it left for maintenance.

\paragraph{200\,W caps.} From about midnight on November~5, GPUs on \texttt{trig0123}, \texttt{0133}, \texttt{0138}, \texttt{0140}, \texttt{0143}, \texttt{0145} and \texttt{0152}
had an enforced limit of 200\,W instead of the default 1,000\,W. According to the administrators, a power glitch at 00:09 caused it, and on the morning of November~6 they cleared it everywhere except on \texttt{trig0145}, which they keep for more tests. Measured under the same matmul load, on every GPU independently: capped GPUs draw \TriCapPowerLo--\TriCapPowerHi\,W, stay at
120\,MHz and reach \TriCapTFlops\,TFLOP/s; GPUs of the normal node \texttt{trig0151} draw \TriCapControlPowerLo--\TriCapControlPowerHi\,W and reach \TriCapControlLo--\TriCapControlHi\,TFLOP/s. \texttt{trig0140}'s
GPU~3 has been back at 1,000\,W and normal since the reboots; its other 7 GPUs have not. One capped node holds back a whole job: a 13-node Qwen run with \texttt{trig0123} took 26.3\,s per step, against about
3\,s without. It was stopped after 8 steps and is used only as evidence for the cap, and it is the only run of ours on a capped node, so no result in this report is affected. The reboots did not clear the caps.

Under the same load after the clearing, the GPUs of the \TriClearedNodes{} cleared nodes draw \TriClearedPowerLo--\TriClearedPowerHi\,W at a 1,000\,W limit and reach \TriClearedTFlopsLo--\TriClearedTFlopsHi\,TFLOP/s, like \texttt{trig0151} beside them (\TriClearedCtlLo--\TriClearedCtlHi). Their health checks pass, and they show the same CPU clock drops as the other nodes. \textbf{One cleared node stays out.} On \texttt{trig0140}, NCCL's all-reduce inside the node fails with default settings, twice out of two, unable to bind NVLink SHARP memory. It works with that feature off, at 647\,GB/s; \texttt{trig0143}, retested beside it, reaches 668\,GB/s with default settings and 647\,GB/s with it off. The same test passed on \texttt{trig0140} on the evening of November~5, on the same boot, so the failure is new, and we do not know why or whether the clearing caused it. Those nodes have been running normally for less than a day, so whether the clearing holds is not yet known either.

\paragraph{CPU clock drops.} In the health check at full GPU power (about 8\,kW per node), mean clocks were \TriCpuBeforeLo--\TriCpuBeforeHi\,MHz (median \TriCpuBeforeMed) on \TriCpuNodesBefore{} nodes before the
kernel change and \TriCpuAfterLo--\TriCpuAfterHi\,MHz (median \TriCpuAfterMed) on \TriCpuNodesAfter{} nodes after it, against 3,200\,MHz at low power. \texttt{trig0114} is the exception, at \TriCpuOddNode\,MHz. In a 1-node Qwen run at 12k tokens per GPU, \texttt{trig0114}
steps take \TriThrottleOddStep\,s against \TriThrottleStep\,s on \texttt{trig0151}, \TriThrottlePct\,\% less (means over the run, repeated runs). This compares two nodes, not a node with and without the drop, so it is an indication of the cost, not a measurement of it.
Pipeline layouts of Qwen (mean 5.8\,kW of GPU power per node) and GLM-5 (4.5\,kW) stay at 3,200\,MHz. Qwen's FSDP and HSDP layouts do not.
\textbf{Not known:} the cause. A node power limit near 10.2\,kW is an inference: the meter's peak matches the GPUs' maximum plus the mean load from everything else, so it is not independent evidence.

\paragraph{Slow nodes.} \texttt{trig0134}, \texttt{0135} and \texttt{0148} spend about half their time below 500\,MHz under Qwen, against about a quarter on \texttt{trig0151}. On 1 node, in the same runs, they are \TriSlowBehindLo--\TriSlowBehindHi\,\% behind \texttt{trig0151}. On 13 nodes, the side-by-side runs with one of them
average \TriSlowThirteenWith\,s per step against \TriSlowThirteenWithout\,s without (+\TriSlowThirteenPct\,\%). That is an association: no test swapped them. All the settings we can read match \texttt{trig0151}, and their CPUs run cooler; the cause is unknown. Runs with these nodes are marked, or left out of scaling comparisons.
A node's low-clock share also depends on the job (\cref{fig:machine_status}, bottom): nodes in a set with a slow node drop less, probably because they wait for it and draw less GPU power (inferred).

\paragraph{Kernel change.} On November~5, from 17:33 to 20:58, the administrators added the boot option \texttt{intel\_pstate=passive}, each node rebooting after its job. The CPU frequency driver changed from \texttt{intel\_pstate} to \texttt{intel\_cpufreq}. Kernel, GPU driver, BIOS, microcode and
power limits are unchanged. \TriKernelReruns{} reruns on the same nodes moved by at most \TriKernelMaxChange\,\%, with repeats before the change at \TriKernelRepeatLo--\TriKernelRepeatHi\,\% (\cref{tab:kernel}); the drops and caps are still there. Results before and after are therefore compared directly, and \texttt{runs.csv} says which is which.
The slow nodes' reruns all moved down, by 0.7--1.0\,\%; two runs a side cannot say whether that is real.

\begin{table}[h]
  \caption{Kernel change: reruns of the same entries on the same nodes. Step times are means over each run's measured steps, one value per run. Repeats: spread of runs on those nodes before the change.}
  \label{tab:kernel}
  \centering\small
  \begin{adjustbox}{max width=\linewidth}
  \IfFileExists{generated/tables/kernel_change.tex}{\input{generated/tables/kernel_change}}{\textit{No data yet.}}
  \end{adjustbox}
\end{table}

\paragraph{Submission.} The site's \texttt{sbatch} wrapper refuses jobs above 25 nodes, and the administrators asked us to submit with the \texttt{.sbatch} binary. A 1-node job submitted that way took 4.767\,s per step, against 4.742 and 4.784\,s through the wrapper. It changed nothing measurable.

\paragraph{Still open.} The administrators' answer on the clock drops, the cap on \texttt{trig0145}, the NCCL failure on \texttt{trig0140}, and the state of \texttt{trig0122} and \texttt{trig0146}.
